% GDZ PPN236018647,native525;HKUST526.
\clearpage\SourceLeaf{525}\ArithmeticHead{496}{L}{TAGEBUCH MIT ERLÄUTERUNGEN.}
\NumberedItem{15}
\indent Coepi de multiplicatoribus (in formis divisorum form[aru]m qu[adraticarum]) connexis cogitare.
\par\noindent\hfill[1796] 22. Iun.\quad G[ottingae]

\medskip\begingroup\small
\indent Vergl. Werke I, S. 476, wo zur Überschrift der Sectio quinta der \textit{Disquisitiones arithmeticae} »de formis aequationibusque indeterminatis secundi gradus« vermerkt ist: Inde a Iun. 22. 1796. Bezüglich der Begriffe »Multiplicatores connexi« und »Forma divisorum« vergl. die oben S. 80 abgedruckte Aufzeichnung, wo im ersten Satze von \textls[120]{Formen der Divisoren} und im art. 1. von \textls[120]{verbundenen}, im art. 4. von \textls[120]{zusammenhängenden} Multiplikatoren die Rede ist, sowie die zugehörigen Bemerkungen S. 83.
\par\noindent\hfill\textsc{Klein. Bachmann.}
\par\endgroup

\NumberedItem{16}
\indent Nova theorematis aurei demonstratio a priori toto coelo diversa eaque haud parum elegans.
\par\noindent\hfill[1796] 27. Iun.

\medskip\begingroup\small
\indent Vergl. Werke I, S. 476, Bemerkung zum art. 262. der \textit{Disquisitiones arithmeticae}, wo indes als Datum für die Auffindung dieses zweiten Beweises des quadratischen Reziprozitätsgesetzes, in Übereinstimmung mit \textsc{Gauss}' eigener Notiz in seinem Handexemplar der \textit{Disquisitiones arithmeticae}, 1796, Juli 27. angegeben ist. Aus der Reihenfolge der Nummern 15, 16, 17 des \textit{Tagebuchs} ergibt sich aber, daß Juli ein Schreibfehler ist. Wegen der Bezeichnung des Reziprozitätsgesetzes der quadratischen Reste als »theorema aureum« vergl. den Artikel 17 des Aufsatzes »Über \textsc{Gauss}' zahlentheoretische Arbeiten«, Werke X 2, S. 44, 45.
\par\noindent\hfill\textsc{Klein. Bachmann.}
\par\endgroup

\NumberedItem{17}
\indent Quaeque partitio numeri \(a\) in tria \(\square\) dat formam in tria \(\square\) separabilem.
\par\noindent\hfill[1796] 3. Iul.

\medskip\begingroup\small
\indent Man vergl. zu dieser und zu der folgenden Nr. 18 die oben S. 78 abgedruckte Aufzeichnung aus dem \textsc{Leiste}, ferner den Artikel 12 des Aufsatzes »Über \textsc{Gauss}' zahlentheoretische Arbeiten«, Werke X 2, S. 29, 30. Unter dem Zeichen \(\square\) ist Quadratzahl bezw. Quadrat einer Linearform, unter \textls[120]{Forma} ohne Zweifel \textls[120]{binäre} quadratische Form zu verstehen; somit weist diese Notiz schon auf die allgemeinen Ergebnisse hin, die im art. 280. der \textit{Disquisitiones arithmeticae}, Werke I, S. 318 dargestellt sind.
\par\noindent\hfill\textsc{Bachmann}.
\par\endgroup
